\section*{الفصل \ref{ch:g7:negatives} --- الأعداد السالبة}

\begin{solution}{exo:g7:negatives:1}
المعاكسات: $-7$؛ \ و $+3.5$؛ \ و $0$ (وهو \omterm{def:g7:negatives:def}{معاكس} نفسه)؛ \ و $-12$. ويكون
عددان متعاكسين حين يكون مجموعهما $0$.
\end{solution}

\begin{solution}{exo:g7:negatives:2}
$-4 < 1$; \quad $-4 > -7$; \quad $-2.5 < -2.4$; \quad $0 > -0.1$.
\end{solution}

\begin{solution}{exo:g7:negatives:3}
$-3.1 < -3 < -0.5 < 0 < 1 < 2.5$.
\end{solution}

\begin{solution}{exo:g7:negatives:4}
$(-5) + (-9) = -14$; \quad
$(-12) + (+7) = -5$; \quad
$(+3.5) + (-1.5) = +2$; \quad
$(-6) + (+6) = 0$.
\end{solution}

\begin{solution}{exo:g7:negatives:5}
$4 - 11 = 4 + (-11) = -7$.

$(-5) - (+3) = (-5) + (-3) = -8$.

$(-7) - (-10) = (-7) + (+10) = +3$.

$0 - (-8) = 0 + (+8) = +8$.
\end{solution}

\begin{solution}{exo:g7:negatives:6}
$(+7) + (-12) + (+3) + (-8) = (+10) + (-20) = -10$ (الموجبات
معًا، والسوالب معًا).

$(-5.5) + (+9) + (-4.5) + (+1) = (-10) + (+10) = 0$.
\end{solution}

\begin{solution}{exo:g7:negatives:7}
$(-3) + 8 - 12 = 5 - 12 = -7$: فهي $-7$ درجة في الليل.
\end{solution}

\begin{solution}{exo:g7:negatives:8}
وضع حرّ؛ والنقطة $E(0, -3)$ تقع على المحور العمودي (فإحداثيها الأول
هو $0$).
\end{solution}

\begin{solution}{exo:g7:negatives:9}
موسكو: $(-4) - (-13) = -4 + 13 = 9$ درجات فرقًا.
وهلسنكي: $2 - (-9) = 11$ درجة.
\end{solution}

\begin{solution}{exo:g7:negatives:10}
$(-12) + 7 - 4 = -5 - 4 = -9$: فهي عند $-9$ م. ولتبلغ السطح
(الارتفاع $0$) عليها أن ترتفع $0 - (-9) = 9$ م.
\end{solution}

\begin{solution}{exo:g7:negatives:11}
$D(-2, -2)$: فإحداثيها الأول كإحداثي $A$، وإحداثيها الثاني كإحداثي
$C$. والأضلاع: $AB = 4$ (من $-2$ إلى $2$) و $BC = 3$ (من $1$ نزولًا إلى
$-2$)، إذن \omterm{def:g6:measure:perimeter}{المحيط} $2 \times (4 + 3) = 14$.
\end{solution}

\begin{solution}{exo:g7:negatives:12}
\omterm{def:g4:geometry:circle}{القطر} الكامل $0, -2, -4$ يعطي \omterm{def:g1:addition:def}{المجموع} السحري
$0 + (-2) + (-4) = -6$. ثم، خلية خلية:
أعلى اليسار $= -6 - (-7) - 0 = 1$؛ ووسط اليمين $= -6 - (-3) - (-2) =
-1$؛ وأسفل اليمين $= -6 - 0 - (-1) = -5$؛ وأسفل الوسط
$= -6 - (-7) - (-2) = 3$.
\begin{center}
\renewcommand{\arraystretch}{1.4}
\begin{tabular}{|c|c|c|}
\hline
$1$ & $-7$ & $0$ \\
\hline
$-3$ & $-2$ & $-1$ \\
\hline
$-4$ & $3$ & $-5$ \\
\hline
\end{tabular}
\end{center}
وكل سطر وعمود \omterm{def:g4:geometry:circle}{وقطر} مجموعه $-6$.
\end{solution}

\begin{solution}{pb:g7:negatives:1}
\textbf{1.} $4\,808 - (-430) = 4\,808 + 430 = 5\,238$ م
(\cref{thm:g7:negatives:sub}: فطرح عدد سالب يعني
الجمع).

\textbf{2.} الرصيد:
\[
(-45) + 120 + (-60) + 15 = 75 + (-60) + 15 = 15 + 15 = 30
\text{ يورو}.
\]
ولإعادة المبلغ الابتدائي $-45$ إلى $0$، أودِع معاكسه:
$45$ يورو.

\textbf{3.} $(-44) - (-100) = (-44) + 100 = 56$: فقد عاش قيصر
$56$ سنة. (ولا فخّ هنا: فالسنتان في الجهة نفسها من
حدّ الحقبتين.)

\textbf{4.} إنها السنة $0$: فالعدّ من $-63$ صعودًا إلى $+14$ يمرّ
بالسنة $0$، لكن تقويم المؤرّخين يقفز مباشرةً من
سنة $1$ قبل الميلاد إلى سنة $1$ بعد الميلاد --- فالسنة $0$ لم توجد قطّ، ومنه فالعدّ
الساذج أطول بسنة واحدة.

\textbf{5.} سنة $63$ قبل الميلاد تصير $-(63 - 1) = -62$. ثم
$14 - (-62) = 14 + 62 = 76$ سنة: أي بالضبط رقم
المؤرّخين. وعلى خطّ الفلكيين المعاد تسميته لا فجوات،
فيعطي الطرح العادي المددَ الحقيقية.

\textbf{6.} $1$ قبل الميلاد $\to 0$؛ \ و $10$ قبل الميلاد $\to -9$؛ \ وبعد الميلاد
$476 \to +476$؛ \ و $753$ قبل الميلاد $\to -752$.

\textbf{7.} $476 - (-752) = 476 + 752 = 1\,228$ سنة.

\textbf{8.} سنة $12$ قبل الميلاد هي $-11$. والظهورات التالية:
$-11 + 76 = 65$، ثم $141$، ثم $217$: أي السنوات $65$ و
$141$ و $217$ بعد الميلاد. وقبلها: $-11 - 76 = -87$، وهي
سنة المؤرّخين $87 + 1 = 88$ قبل الميلاد.

\textbf{9.} سنة $5$ قبل الميلاد هي $-4$، إذن $5 - (-4) = 9$ سنوات --- لا
$10$. فسنة \omterm{def:g1:counting:zero}{الصفر} المفقودة تضرب كلما عبرت مدّة
حدّ الحقبتين.

\textbf{10.} سنة $3$ قبل الميلاد هي $-2$؛ وبعد عشر سنوات $-2 + 10 = +8$:
فقد بلغ الطفل العاشرة سنة $8$ بعد الميلاد.

\textbf{11.} بين $-7$ و $-2$: $(-2) - (-7) = 5$. وبين
$-3.5$ و $4.5$: $4.5 - (-3.5) = 8$.

\textbf{12.} بتجميع المتعاكسات والأصدقاء:
\[
3 + (-5) + 2 + (-1) + 4
= \bigl(3 + 2\bigr) + (-5) + (-1) + 4
= 0 + (-1) + 4 = +3 :
\]
فينتهي المصعد عند الطابق $3$. والترتيب لا يهمّ: \omterm{def:g1:addition:def}{فمجموع}
الأعداد النسبية يمكن إعادة تنظيمه بحرّية (فكل حركة تُمشى
على \omterm{def:g6:lines:objects}{المستقيم} نفسه، والمشية الكلية هي نفسها).

\textbf{13.} التغيّر الصافي:
$2 + (-4) + 1 + 0 + (-3) + 5 + (-2) = -1$ درجة. واِنطلاقًا من
$-2$: تكون درجة الحرارة النهائية $-2 + (-1) = -3$ درجة. وعمومًا،
القيمة النهائية هي القيمة الابتدائية زائد \omterm{def:g1:addition:def}{مجموع} كل
التغيّرات --- ولا حاجة إلى تتبّع الصعود والنزول واحدًا واحدًا.

\textbf{14.} العودة إلى البداية تعني تغيّرًا صافيًا $0$: فالصعود
والنزول يجب أن يتلاشيا بالضبط. \omterm{def:g1:addition:def}{ومجموع} التغيّرات الخمسة المعلومة
$4 + (-7) + 2 + 3 + (-5) = -3$، إذن التغيّر السادس هو
\omterm{def:g7:negatives:def}{المعاكس}: $+3$ (أي ثلاثمئة متر صعودًا).

\textbf{15.} عند الفلكيين: سنة $509$ قبل الميلاد $\to -508$، و
$476 - (-508) = 984$ سنة. وحساب المؤرّخ $509 + 476 = 985$
يعدّ سنة زائدة --- ومرة أخرى إنها سنة \omterm{def:g1:counting:zero}{الصفر} الشبح،
المدرجة صامتةً في الجمع والغائبة عن التاريخ.
(وللتحقّق، الحقبتان: الجمهورية من $-508$ إلى $-26$، أي $482$
سنة؛ والإمبراطورية من $-26$ إلى $+476$، أي $502$ سنة؛
و $482 + 502 = 984$.)
\end{solution}
